\ifdefined\idpCombinedMode
\else
\documentclass{idpbeamer}
\fi


\lectureinfo{POO Java}{Aula 02\\Plataforma Java}{Prof. Dr. Nome do Professor}

\startLectureFile

\section{Plataforma}

\begin{frame}{Plataforma}
A plataforma Java reúne linguagem, bibliotecas, ferramentas e um ambiente de execução portável.


\begin{itemize}
  \item JDK: ferramentas de desenvolvimento.
  \item JRE: ambiente de execução.
  \item JVM: máquina virtual que executa bytecode.
\end{itemize}
\end{frame}

\section{Fluxo}

\begin{frame}{Fluxo de Execução}
\begin{center}
\Large
\texttt{.java} $\rightarrow$ \texttt{.class} $\rightarrow$ JVM
\end{center}

O compilador transforma código-fonte em bytecode, e a JVM executa esse bytecode no sistema operacional disponível.
\end{frame}

\finishLectureFile
